\section{Finite-mode numerical realization}\label{sec:numerics}

The preceding continuum solutions admit a common numerical treatment: expand each cut region in modes of its decoupled spatial operator, retain finitely many coefficients, add the interface quadratic form before diagonalization, and compare the resulting low spectrum with an independent continuum secular equation. This section applies that procedure only to the models treated in the present article. The cutoff label $N$ means that the mode indices $0,\ldots,N$ are retained unless stated otherwise.

\subsection{Truncation, diagonalization, and response matching}\label{subsec:numerics-method}

Let $Q$ and $P$ collect the retained canonical coefficients. Their quadratic Hamiltonian has the form

\begin{align}
H_N
=\frac12P^{\mathrm T}G_N^{-1}P
+\frac12Q^{\mathrm T}K_NQ.
\end{align}

If the interface action contains no time derivatives, $G_N=I$ and the squared frequencies are the eigenvalues of $K_N$. If the interface degrees of freedom contribute to the CPS form, one instead solves

\begin{align}
K_Nv_j=\omega_j^2G_Nv_j,
\qquad
v_i^{\mathrm T}G_Nv_j=\delta_{ij}.
\end{align}

The second normalization is the finite-dimensional form of CPS normalization. In particular, replacing it by the Euclidean norm when $G_N\neq I$ changes the physical mode normalization.

For a scalar stiffness defect, write

\begin{align}
K_N=D_N+C_N^{\mathrm T}A_NC_N,
\end{align}

where $D_N$ is the diagonal decoupled stiffness, $C_N$ evaluates the retained fields on the cut, and $A_N$ is the bare boundary matrix. Define the truncated boundary resolvent

\begin{align}
R_N(z)=C_N(D_N-zI)^{-1}C_N^{\mathrm T}.
\end{align}

Away from a free pole, the matrix determinant lemma gives

\begin{align}
\frac{\det(K_N-zI)}{\det(D_N-zI)}
=\det(A_N)\det\!\left(A_N^{-1}+R_N(z)\right).
\end{align}

Let $R(z)$ be the continuum response and $T_N(z)=R(z)-R_N(z)$ the omitted tail. Matching at a reference value $z_*$ is achieved by

\begin{align}
A_N=\left(A^{-1}+T_N(z_*)\right)^{-1},
\end{align}

because then

\begin{align}
A_N^{-1}+R_N(z_*)=A^{-1}+R(z_*).
\end{align}

We use $z_*=0$ in the spatial spectral variable. The residual secular-matrix error is exactly $T_N(z)-T_N(0)$. Its conversion into an eigenvalue error assumes an isolated, nondegenerate continuum root; consequently the observed cutoff power is a low-energy, fixed-window statement and is not uniform up to the ultraviolet cutoff. Eigenvalues at free poles are always checked in the original matrix problem rather than in the divided determinant.

Each calculation below performs four independent checks: positivity of the finite quadratic data, a matrix or symmetry identity, comparison with an analytic spectrum or a finer regulator, and convergence in a fixed spectral window. The accompanying Mathematica files are \texttt{finite\_truncation.wl} and the scripts under \texttt{numerics/}.

\subsection{Interval scalars}\label{subsec:numerics-interval}

For the Neumann model of Section $\ref{sec:interval-scalar}$, use the normalized half-interval bases

\begin{align}
u_{1,0}=\frac1{\sqrt L},
\qquad
u_{1,n}=\sqrt{\frac2L}\cos\!\left(\frac{n\pi(x+L)}L\right),
\end{align}

\begin{align}
u_{2,0}=\frac1{\sqrt L},
\qquad
u_{2,n}=\sqrt{\frac2L}\cos\!\left(\frac{n\pi(L-x)}L\right),
\qquad n\geq1.
\end{align}

Let $c$ be the vector of left-minus-right cut values. Then

\begin{align}
K_N
=\operatorname{diag}\!\left(
m^2+\frac{n^2\pi^2}{L^2},
m^2+\frac{n^2\pi^2}{L^2}
\right)_{n=0}^{N}
+g_Ncc^{\mathrm T}.
\end{align}

The common sector lies in the kernel of $c^{\mathrm T}$ and is reproduced exactly at every cutoff. The relative sector converges to the roots of $k\tan(kL)=2g$. For perfect gluing, direct projection onto $c^{\mathrm T}Q=0$ removes one finite-dimensional direction. The leading response-matched penalty

\begin{align}
g_N^{(\infty)}=\frac{\pi^2N}{4L}
\end{align}

instead retains the direction while compensating the leading omitted boundary response. With $L=m=1$, the maximum relative error among the first eight full-interval frequencies is

\begin{align}
\begin{array}{c|cc}
N & \text{hard projection} & \text{response matched} \\ \hline
8   & 2.5151\times10^{-2} & 1.2114\times10^{-3} \\
16  & 1.2445\times10^{-2} & 3.3827\times10^{-4} \\
32  & 6.2371\times10^{-3} & 8.9449\times10^{-5} \\
64  & 3.1276\times10^{-3} & 2.3005\times10^{-5} \\
128 & 1.5668\times10^{-3} & 5.8738\times10^{-6}
\end{array}
\end{align}

Thus the hard projection is $O(N^{-1})$ in this benchmark, while the matched penalty is $O(N^{-2})$ on the displayed fixed window.

The Dirichlet interpolation of Appendix $\ref{app:dirichlet}$ requires a different off-shell basis. Modes that obey strict Dirichlet data at the cut have zero cut trace and therefore cannot represent the boundary penalty. We instead impose Dirichlet data only at the physical endpoint and use cut-Neumann modes

\begin{align}
k_n^{(0)}=\frac{(n+\frac12)\pi}{L},
\qquad
b_n=(-1)^n\sqrt{\frac2L}.
\end{align}

In the common and relative sectors the boundary stiffnesses are

\begin{align}
\alpha_+=\frac\mu\kappa,
\qquad
\alpha_-=\mu\left(\frac1\kappa+2\kappa\right),
\end{align}

and each finite matrix is

\begin{align}
(K_{\sigma,N})_{nr}
=\left[m^2+(k_n^{(0)})^2\right]\delta_{nr}
+\alpha_{\sigma,N}b_nb_r.
\end{align}

The exact comparison roots obey $k\cot(kL)=-\alpha_\sigma$. Since

\begin{align}
T_N(0)
=\sum_{n=N+1}^{\infty}\frac{b_n^2}{(k_n^{(0)})^2}
=\frac{2L}{\pi^2}\psi_1\!\left(N+\frac32\right),
\end{align}

where $\psi_1$ is the trigamma function. The two response-matched coefficients satisfy

\begin{align}
\frac1{\alpha_{\sigma,N}}
=\frac1{\alpha_\sigma}+T_N(0),
\qquad \sigma=\pm.
\end{align}

For $L=m=\mu=1$, the maximum relative error among the first eight frequencies falls from $5.19\times10^{-3}$ to $1.17\times10^{-5}$ at $N=16$, and from $1.34\times10^{-3}$ to $2.04\times10^{-7}$ at $N=64$, for $\kappa=1$. The same test was performed at $\kappa=0.1$ and $5$.

This example also exposes a structural consequence of the cutoff. Writing the regulated boundary form as

\begin{align}
d_N(q_1^2+q_2^2)+h_N(q_1-q_2)^2,
\end{align}

the matched values are

\begin{align}
d_N=\alpha_{+,N},
\qquad
h_N=\frac{\alpha_{-,N}-\alpha_{+,N}}2.
\end{align}

In general $d_Nh_N\neq\mu^2$. The continuum one-parameter curve $d=\mu/\kappa$, $h=\mu\kappa$ is therefore not closed under low-energy cutoff matching: both symmetry-allowed boundary coefficients must run independently if both parity sectors are to be accelerated.

More explicitly, for $\alpha_+=d$, $\alpha_-=d+2h$, and $dh=\mu^2$,

\begin{align}
d_Nh_N
=\frac{\mu^2}{
\left(1+\alpha_-T_N(0)\right)
\left(1+\alpha_+T_N(0)\right)^2}.
\end{align}

This exact identity makes the failure of a single bare $\kappa_N$ quantitative.

\subsection{Scalar geometries}\label{subsec:numerics-scalar-geometries}

For global $\mathrm{AdS}_2$, restrict the normalized even global modes to a half-space. Their free frequencies and origin traces are

\begin{align}
\omega_r^{(0)}=\Delta+2r,
\qquad
b_r=\sqrt2\,
\frac{C_{2r}^{\Delta}(0)}{\sqrt{h_{2r}^{(\Delta)}}}.
\end{align}

The reflection-even tower remains exact, while the affected reflection-odd matrix is

\begin{align}
(K_N)_{rs}
=(\Delta+2r)^2\delta_{rs}+2g_Nb_rb_s.
\end{align}

The continuum response in Section $\ref{subsec:global-ads2}$ implies

\begin{align}
\sum_{r=0}^{\infty}
\frac{b_r^2}{(\Delta+2r)^2-\omega^2}
=-\frac1{D_\Delta(\omega)},
\end{align}

so, with $T_N^{\mathrm{AdS}_2}(0)$ the omitted part of this sum,

\begin{align}
\frac1{g_N}=\frac1g+2T_N^{\mathrm{AdS}_2}(0).
\end{align}

At $\Delta=2$ and $g=1$, comparison with the exact roots $D_\Delta(\omega)=2g$ gives

\begin{align}
\begin{array}{c|cc}
N & \text{direct error} & \text{response-matched error} \\ \hline
8  & 7.5922\times10^{-3} & 1.0363\times10^{-4} \\
16 & 4.0363\times10^{-3} & 1.3207\times10^{-5} \\
32 & 2.0864\times10^{-3} & 1.7801\times10^{-6} \\
64 & 1.0615\times10^{-3} & 2.3468\times10^{-7}
\end{array}
\end{align}

where the error is the maximum relative error among the first six affected frequencies. This confirms that response matching is not tied to trigonometric flat-space modes.

For the four-quadrant scalar, use a product Neumann basis on each quadrant. The matrix dimension is $4(N+1)^2$. A vertical half-axis penalty is rank one in the retained $x$ coefficients for each fixed $y$ mode; a horizontal penalty is the corresponding rank-one update in $y$. All four updates are positive for nonnegative couplings.

When $g_x^-=g_x^+$ and $g_y^-=g_y^+$, the assembled matrix equals the Kronecker sum of the two one-dimensional interval matrices. At $N=4$ the spectral residual of this identity is $6.3\times10^{-13}$. For four unrelated couplings, exchanging $x$ and $y$ together with the coupling pairs preserves the spectrum to $1.2\times10^{-12}$, and the constant mode remains exact.

For a local half-axis coupling $g_e$, the scalar response prescription gives

\begin{align}
\frac1{g_{e,N}}
=\frac1{g_e}
+\frac{4L}{\pi^2}\psi_1(N+1).
\end{align}

For

\begin{align}
(g_x^-,g_x^+,g_y^-,g_y^+)=(0.4,1.1,0.7,2.0),
\qquad L=m=1,
\end{align}

the maximum absolute error among the first twelve frequencies, measured against a matched $N=36$ reference, is

\begin{align}
\begin{array}{c|cc}
N & \text{direct error} & \text{response-matched error} \\ \hline
4  & 8.6549\times10^{-2} & 1.9633\times10^{-3} \\
8  & 4.3075\times10^{-2} & 2.8252\times10^{-4} \\
16 & 2.1578\times10^{-2} & 3.5314\times10^{-5} \\
24 & 1.4403\times10^{-2} & 8.2818\times10^{-6}
\end{array}
\end{align}

The separable calculation is an exact finite-matrix identity. The nonseparable calculation is only refinement evidence because the $N=36$ reference is not an analytic continuum spectrum.

\subsection{Maxwell kinetic defect}\label{subsec:numerics-maxwell}

After the reduction of Section $\ref{subsec:maxwell-dual-scalar}$, decompose along the circle and then diagonalize the $x$ dependence. The unaffected reflection-odd tower has zero cut trace. For the reflection-even tower, use the outer-Dirichlet/cut-Neumann basis

\begin{align}
k_n^{(0)}=\frac{(n+\frac12)\pi}{L},
\qquad
b_n=\frac{(-1)^n}{\sqrt L}.
\end{align}

The bulk gradient term supplies the stiffness matrix, whereas the cut term in the reduced CPS form supplies a rank-one kinetic defect:

\begin{align}
K_N=\operatorname{diag}\!\left((k_n^{(0)})^2\right),
\qquad
G_N=I+\frac1g bb^{\mathrm T}.
\end{align}

For $g>0$, $G_N$ is positive. The generalized eigenproblem

\begin{align}
K_Nv=k^2G_Nv
\end{align}

has the finite secular equation

\begin{align}
1-\frac{k^2}{g}
\sum_{n=0}^{N}\frac{b_n^2}{(k_n^{(0)})^2-k^2}=0.
\end{align}

The continuum response identity

\begin{align}
\sum_{n=0}^{\infty}\frac{b_n^2}{(k_n^{(0)})^2-k^2}
=\frac{\tan(kL)}{2k}
\end{align}

reproduces $k\tan(kL)=2g$. The circle momentum is then restored algebraically by $\omega_{m,j}^2=q_m^2+k_j^2$; it does not require a second diagonalization.

For $L=g=1$, the maximum absolute error among the first eight affected roots is

\begin{align}
\begin{array}{c|c}
N & \max_{1\leq j\leq8}|k_{j,N}-k_j| \\ \hline
8  & 5.7591\times10^{-1} \\
16 & 2.6686\times10^{-1} \\
32 & 1.3439\times10^{-1} \\
64 & 6.7957\times10^{-2}
\end{array}
\end{align}

At $N=24$, the generalized eigen-equation residual is $7.0\times10^{-37}$, the $G_N$-orthogonality residual is $1.0\times10^{-38}$, and the smallest eigenvalue of $G_N$ is one. These checks directly test the interface contribution to the CPS norm.

This kinetic defect does not admit the scalar zero-energy matching. If

\begin{align}
T_N(k)
=\frac{\tan(kL)}{2k}
-\sum_{n=0}^{N}\frac{b_n^2}{(k_n^{(0)})^2-k^2},
\end{align}

then exact matching at a selected root requires

\begin{align}
g_N(k)=g-k^2T_N(k).
\end{align}

The correction vanishes at $k=0$ and depends on the selected energy. Hence one constant bare coupling cannot accelerate an entire finite window. A local realization would require higher-derivative interface terms, or equivalently an auxiliary-mode completion. The common holonomy is a separate exact canonical pair; the compact magnetic-flux components in Appendix $\ref{subsec:dirichlet-maxwell}$ are discrete sectors and are not produced by oscillator diagonalization.

The Maxwell Dirichlet interpolation of Appendix $\ref{subsec:dirichlet-maxwell}$ has the same kinetic-defect structure in two boundary sectors. Diagonalizing $M_\kappa$ gives

\begin{align}
\alpha_+=\frac\mu\kappa,
\qquad
\alpha_-=\mu\left(\frac1\kappa+2\kappa\right).
\end{align}

Using half-interval outer-Dirichlet/cut-Neumann modes with

\begin{align}
k_n^{(0)}=\frac{(n+\frac12)\pi}{L},
\qquad
b_n=(-1)^n\sqrt{\frac2L},
\end{align}

the two pencils are

\begin{align}
K_N=\operatorname{diag}\!\left((k_n^{(0)})^2\right),
\qquad
G_{\sigma,N}=I+\frac1{\alpha_\sigma}bb^{\mathrm T}.
\end{align}

Their continuum roots obey

\begin{align}
k\tan(kL)=\alpha_\sigma.
\end{align}

As $\kappa\to0$, both sectors approach the outer-Dirichlet/cut-Neumann half-interval spectrum. As $\kappa\to\infty$, the common and relative roots approach $n\pi/L$ and $(n+\frac12)\pi/L$, respectively. Apart from the zero root that becomes the common holonomy, their union is the full-interval Dirichlet photon spectrum.

For $L=\mu=1$, the maximum absolute error among the first eight combined roots is

\begin{align}
\begin{array}{c|ccc}
N & \kappa=0.1 & \kappa=1 & \kappa=5 \\ \hline
8  & 1.0595\times10^{-1} & 2.1519\times10^{-1} & 2.2117\times10^{-1} \\
16 & 5.6905\times10^{-2} & 1.1188\times10^{-1} & 1.1408\times10^{-1} \\
32 & 2.9775\times10^{-2} & 5.7358\times10^{-2} & 5.8228\times10^{-2} \\
64 & 1.5272\times10^{-2} & 2.9082\times10^{-2} & 2.9455\times10^{-2}
\end{array}
\end{align}

Here the omitted response is $T_N(k)=\tan(kL)/k-R_N(k^2)$, and selected-root matching again requires $\alpha_{\sigma,N}(k)=\alpha_\sigma-k^2T_N(k)$. Both CPS metrics are positive and were normalized in the generalized eigenbasis. The discrete compact-$U(1)$ sectors remain outside this oscillator calculation.

The strong-coupling endpoint also depends on the scaling used to take the limit. Fixed unscaled gauge-invariant data give the pure uncut Maxwell phase space after pullback and quotient. A family whose common boundary component scales as $\kappa^{1/2}$ instead retains the additional massless interface scalar described by the full continuum analysis. The direct truncation above follows the ordinary oscillator roots; it does not identify these two limiting phase spaces.

\subsection{Scope of the numerical evidence}\label{subsec:numerics-scope}

The calculations support four analytic conclusions. First, the continuum boundary response, not a regulator-dependent penalty coefficient, is the quantity that survives comparison across flat intervals, global $\mathrm{AdS}_2$, and interface networks. Second, a continuum one-parameter interface family need not be closed under truncation: the Dirichlet interpolation generates independent running coefficients in its common and relative sectors. Third, an interface contribution to the CPS form changes the finite spectral problem itself, as demonstrated by both Maxwell kinetic pencils. Fourth, exact protected sectors---the scalar common tower, the four-quadrant constant mode, the Maxwell odd tower, and the common holonomy---provide sensitive checks of signs and interface assembly.

What is established here is a collection of exact finite-matrix identities and fixed-window spectral convergence tests. These results do not establish convergence of unsmeared fields or correlators, equivalence of continuum Fock representations, or an interacting continuum gluing construction. The nonseparable four-quadrant spectrum has only a refinement reference, and an energy-independent accelerated Maxwell regulator has not been constructed.
